% 数理物理1 数理物理1-9 問題4（060-real-n4。当時は 数理物理1-7 問題7・shou07-l2-real-n4） 解いた日 2026-10-12
%% kotae
\toi{1}{○}{記述}
（ページの上に $\sin a\cos b$ のメモ）
\begin{align*}
  &\sin\frac{2\pi\mu n}{N}\cos\frac{2\pi\nu n}{N}\\
  &\qquad=\frac{1}{2}\Bigl(\sin\frac{2\pi(\mu+\nu)n}{N}+\sin\frac{2\pi(\mu-\nu)n}{N}\Bigr)\\
  &\qquad=\frac{1}{2}\Im\Bigl(\ee^{\ii\frac{2\pi(\mu+\nu)n}{N}}+\ee^{\ii\frac{2\pi(\mu-\nu)n}{N}}\Bigr)
\end{align*}
$z_1=\ee^{\ii\frac{2\pi(\mu+\nu)}{N}}$，$z_2=\ee^{\ii\frac{2\pi(\mu-\nu)}{N}}$ とする。
\begin{align*}
  \text{与式}&=\sum_{n=1}^{N}\frac{1}{2}\Im\bigl(z_1^{n}+z_2^{n}\bigr)\\
  &=0
\end{align*}
（脇に）$\mu+\nu$ や $\mu-\nu$ が $N$ の倍数でも倍数でなくても虚部は 0。
\naze{虚部が 0 になる理由（和の値）が無い。$S_p=\sum_{n=1}^{N}z^{n}$ は，$p$ が $N$ の倍数なら各項が 1 で $N$，そうでなければ $z\neq1$，$z^{N}=1$ より $z(1-z^{N})/(1-z)=0$ で，どちらも実数なので $\Im S_p=0$ である。}
\toi{2}{◎}{}
\begin{align*}
  A_{\mu}&=\frac{1}{4}\sum_{n=1}^{4}\cos\frac{\pi\mu n}{2}f_n\\
  A_1&=\frac{1}{4}\Bigl(\cos\frac{\pi}{2}+\cos\pi\Bigr)=-\frac{1}{4}\\
  A_2&=\frac{1}{4}(\cos\pi+\cos2\pi)=0\\
  A_3&=\frac{1}{4}\Bigl(\cos\frac{3\pi}{2}+\cos3\pi\Bigr)=-\frac{1}{4}\\
  A_4&=\frac{1}{4}(\cos2\pi+\cos4\pi)=\frac{1}{2}
\end{align*}
\begin{align*}
  B_{\mu}&=\frac{1}{4}\sum_{n=1}^{4}\sin\frac{\pi\mu n}{2}f_n\\
  B_1&=\frac{1}{4}\Bigl(\sin\frac{\pi}{2}+\sin\pi\Bigr)=\frac{1}{4}\\
  B_2&=\frac{1}{4}(\sin\pi+\sin2\pi)=0\\
  B_3&=\frac{1}{4}\Bigl(\sin\frac{3\pi}{2}+\sin3\pi\Bigr)=-\frac{1}{4}\\
  B_4&=\frac{1}{4}(\sin2\pi+\sin4\pi)=0
\end{align*}
\toi{3}{○}{記述}
\[ \hat{f}_{\mu}=\frac{1}{2}\sum_{n}\bigl(\ee^{-\ii\frac{\pi\mu n}{2}}\bigr)_{\kakimachigai{n\mu}{\mu n}}f_n \]
$\ee^{\ii\frac{\pi}{2}}=z$ とする。
\begin{align*}
  2U^{\dagger}&=\begin{pmatrix}z^{-1}&z^{-2}&z^{-3}&z^{-4}\\z^{-2}&z^{-4}&z^{-6}&z^{-8}\\z^{-3}&z^{-6}&z^{-9}&z^{-12}\\z^{-4}&z^{-8}&z^{-12}&z^{-16}\end{pmatrix}\\
  &=\begin{pmatrix}z^{3}&z^{2}&z&1\\z^{2}&1&z^{2}&1\\z&z^{2}&z^{3}&1\\1&1&1&1\end{pmatrix}
\end{align*}
\[
  2U^{\dagger}\vect{f}=\begin{pmatrix}z^{3}+z^{2}\\z^{2}+1\\z+z^{2}\\2\end{pmatrix}=\begin{pmatrix}-\ii-1\\0\\\ii-1\\2\end{pmatrix}
\]
\begin{align*}
  \hat{f}_{\mu}&=\Bigl(-\frac{1}{2}-\frac{\ii}{2},\ 0,\ -\frac{1}{2}+\frac{\ii}{2},\ 1\Bigr)\\
  \frac{1}{2}\Re\hat{f}_{\mu}&=\Bigl(-\frac{1}{4},\ 0,\ -\frac{1}{4},\ \frac{1}{2}\Bigr)\\
  -\frac{1}{2}\Im\hat{f}_{\mu}&=\Bigl(\frac{1}{4},\ 0,\ -\frac{1}{4},\ 0\Bigr)
\end{align*}
\nuke{結論の1行}
\naze{値は (2) の $A_{\mu}$，$B_{\mu}$ と一致しているが，「よって $\Re\hat{f}_{\mu}=2A_{\mu}=\sqrt{N}A_{\mu}$，$\Im\hat{f}_{\mu}=-2B_{\mu}=-\sqrt{N}B_{\mu}$ が成り立つ」という確かめの結論が無い。「確かめよ」の問いでは，比べた結果を1行で書く。}
\toi{4}{◎}{}
\begin{align*}
  \kakimachigai{\abs{\hat{f}_{\mu}}^{2}}{\textstyle\sum_{\mu}\abs{\hat{f}_{\mu}}^{2}}&=\frac{1}{2}+0+\frac{1}{2}+1=2\\
  \kakimachigai{\abs{f_n}^{2}}{\abs{\vect{f}}^{2}}&=1+1+0+0=2
\end{align*}
（0 になる項は斜線で消して代入する）
\begin{align*}
  f_1&=B_1\sin\frac{\pi}{2}+B_3\sin\frac{3\pi}{2}+A_4\cos2\pi\\
  &=\frac{1}{4}+\frac{1}{4}+\frac{1}{2}=1\\
  f_2&=A_1\cos\pi+A_2\cos2\pi+A_3\cos3\pi+A_4\cos4\pi\\
  &=\frac{1}{4}+0+\frac{1}{4}+\frac{1}{2}=1\\
  f_3&=B_1\sin\frac{3\pi}{2}+A_2\cos3\pi+B_3\sin\frac{9\pi}{2}+A_4\cos6\pi\\
  &=-\frac{1}{4}+0-\frac{1}{4}+\frac{1}{2}=0\\
  f_4&=A_1\cos2\pi+A_2\cos4\pi+A_3\cos6\pi+A_4\cos8\pi\\
  &=-\frac{1}{4}+0-\frac{1}{4}+\frac{1}{2}=0
\end{align*}
%% kizuki
\item \mae{数理物理1-7（組み替え前）問題6 (1)・10-03} 「和が 0（虚部が 0）」と言うときは，1の $N$ 乗根の和 $S_p$ の値を書く。$p$ が $N$ の倍数なら $z=1$ で $S_p=N$，そうでなければ $z\neq1$，$z^{N}=1$ で $S_p=z(1-z^{N})/(1-z)=0$。どちらも実数だから $\Im S_p=0$ と言える。等比級数の式を書きかけて消していたので，最後まで書けば足りた。
\item \mae{数理物理1-7（組み替え前）問題7 (3)・10-03} 前回は $\hat{f}_{\mu}$ の値が無く，今回は値は出たが「$\Re\hat{f}_{\mu}=\sqrt{N}A_{\mu}$，$\Im\hat{f}_{\mu}=-\sqrt{N}B_{\mu}$ が成り立つ」の結論が無い。「確かめよ」は比べた結果を1行で書いて終える。前回の (4) の Parseval（絶対値の2乗と因子 4）は今回正しくできた。
